%; whizzy paragraph -pdf xpdf -latex ./whizzypdfptex.sh
%; whizzy-paragraph "^\\\\begin{frame}\\|\\\\emtext"
% latex beamer presentation.
% platex, latex-beamer でコンパイルすることを想定。 

%     Tokyo Debian Meeting resources
%     Copyright (C) 2012 Junichi Uekawa

%     This program is free software; you can redistribute it and/or modify
%     it under the terms of the GNU General Public License as published by
%     the Free Software Foundation; either version 2 of the License, or
%     (at your option) any later version.

%     This program is distributed in the hope that it will be useful,
%     but WITHOUT ANY WARRANTY; without even the implied warreanty of
%     MERCHANTABILITY or FITNESS FOR A PARTICULAR PURPOSE.  See the
%     GNU General Public License for more details.

%     You should have received a copy of the GNU General Public License
%     along with this program; if not, write to the Free Software
%     Foundation, Inc., 51 Franklin St, Fifth Floor, Boston, MA  02110-1301 USA

\documentclass[cjk,dvipdfmx,12pt]{beamer}
\usetheme{Tokyo}
\usepackage{monthlypresentation}

%  preview (shell-command (concat "evince " (replace-regexp-in-string "tex$" "pdf"(buffer-file-name)) "&")) 
%  presentation (shell-command (concat "xpdf -fullscreen " (replace-regexp-in-string "tex$" "pdf"(buffer-file-name)) "&"))
%  presentation (shell-command (concat "evince " (replace-regexp-in-string "tex$" "pdf"(buffer-file-name)) "&"))

%http://www.naney.org/diki/dk/hyperref.html
%日本語EUC系環境の時
\AtBeginDvi{\special{pdf:tounicode EUC-UCS2}}
%シフトJIS系環境の時
%\AtBeginDvi{\special{pdf:tounicode 90ms-RKSJ-UCS2}}

\newenvironment{commandlinesmall}%
{\VerbatimEnvironment
  \begin{Sbox}\begin{minipage}{1.0\hsize}\begin{fontsize}{8}{8} \begin{BVerbatim}}%
{\end{BVerbatim}\end{fontsize}\end{minipage}\end{Sbox}
  \setlength{\fboxsep}{8pt}
% start on a new paragraph

\vspace{6pt}% skip before
\fcolorbox{dancerdarkblue}{dancerlightblue}{\TheSbox}

\vspace{6pt}% skip after
}
%end of commandlinesmall

\title{東京エリアDebian勉強会}
\subtitle{第126回 2015年5月度}
\author{野島貴英}
\date{2015年5月23日}
\logo{\includegraphics[width=8cm]{image200607/openlogo-light.eps}}

\begin{document}

\begin{frame}
\titlepage{}
\end{frame}

\begin{frame}{設営準備にご協力ください。}
会場設営よろしくおねがいします。
\end{frame}

\begin{frame}{Agenda}
 \begin{minipage}[t]{0.45\hsize}
  \begin{itemize}
   \item 注意事項
	 \begin{itemize}
	  \item 写真はセミナールーム内のみ可です。
          \item 出入りは自由でないので、もし外出したい方は、野島まで一声くださいませ。
	 \end{itemize}
   \item 事前課題発表
  \end{itemize}
 \end{minipage} 
 \begin{minipage}[t]{0.45\hsize}
  \begin{itemize}
   \item 最近あったDebian関連のイベント報告
	 \begin{itemize}
	 \item 第125回 東京エリアDebian勉強会
	 \end{itemize}
   \item Debian Trivia Quiz
   \item 自然言語処理チームとパッケージ
   \item 今後のイベント
   \item 今日の宴会場所
  \end{itemize}
 \end{minipage}
\end{frame}

\section{事前課題}
\emtext{事前課題}
{\footnotesize
\begin{prework}{ $BLnEg(B }
  \begin{enumerate}
  \item Q.hack time$B$K2?$r$7$^$9$+!)(B\\
    A. Nook HD+$B$r$=$m$=$m(Bdebian$B$K!#(B
  \item ($B%*%W%7%g%s(B)Q.$B2?$K$D$$$FJ9$-$?$$!?;22C<T$HOC$r$7$?$$$G$9$+!)(B\\
    A. $B%=%U%H%&%'%"<+M3$r<i$jDL$9;v$KBP$7$F!"2?$,=PMh$k$+$K$D$$$F!#(B
  \end{enumerate}
\end{prework}

\begin{prework}{ wskoka }
  \begin{enumerate}
  \item Q.hack time$B$K2?$r$7$^$9$+!)(B\\
    A. MIPS Debian $B$N0\?"(B
  \item ($B%*%W%7%g%s(B)Q.$B$I$3$G:#2s$NJY6/2q$N3+:E$rCN$j$^$7$?$+!)(B\\
    A. $B$=$NB>(B
  \end{enumerate}
\end{prework}

\begin{prework}{ koedoyoshida  }
  \begin{enumerate}
  \item Q.hack time$B$K2?$r$7$^$9$+!)(B\\
    A. DDTSS,PyconJP$B4XO"(B
  \item ($B%*%W%7%g%s(B)Q.$B$I$3$G:#2s$NJY6/2q$N3+:E$rCN$j$^$7$?$+!)(B\\
    A. $BM'C#$dCN$j9g$$$+$iD>@\(B
  \end{enumerate}
  DDTSS: http://ddtp.debian.net/ddtss/index.cgi/ja
\end{prework}

\begin{prework}{ dictoss }
  \begin{enumerate}
  \item Q.hack time$B$K2?$r$7$^$9$+!)(B\\
    A. kFreeBSD$B$N(BIPsec$B4XO"$rD4$Y$k(B
  \item ($B%*%W%7%g%s(B)Q.$B$I$3$G:#2s$NJY6/2q$N3+:E$rCN$j$^$7$?$+!)(B\\
    A. Debian JP$B$N%a!<%j%s%0%j%9%H(B
  \end{enumerate}
\end{prework}

\begin{prework}{ yy\_y\_ja\_jp }
  \begin{enumerate}
  \item Q.hack time$B$K2?$r$7$^$9$+!)(B\\
    A. DDTSS 
  \item ($B%*%W%7%g%s(B)Q.$B$I$3$G:#2s$NJY6/2q$N3+:E$rCN$j$^$7$?$+!)(B\\
    A. $B$=$NB>(B
  \item ($B%*%W%7%g%s(B)Q.$B2?$K$D$$$FJ9$-$?$$!?;22C<T$HOC$r$7$?$$$G$9$+!)(B\\
    A. DDTSS
  \end{enumerate}
\end{prework}

}

\section{イベント報告}
\emtext{イベント報告}

\begin{frame}{第125回東京エリアDebian勉強会}

\begin{itemize}
\item 場所はスクウェア・エニックスさんのセミナルームをお借りしての開催でした。
\item 参加者は13名でした。
\item セミナ内容はまえだこうへいさんによる「.debからPythonパッケージへの変遷」でした。
\item 残りの時間でhack timeを行い、成果発表をしました。
\item 宴会の代わりに、「祥龍房 新宿イーストサイドスクエア店」で夕食会をやりました。
\end{itemize} 
  
\end{frame}

\begin{frame}{第125回東京エリアDebian勉強会(つづき)}

 セミナはまえださんにより、PythonパッケージのローカルDebian CIに関する発表と、Pythonに関するDebian公式パッケージの考察が行われました。大統一Debian勉強会から、毎年発表のはや３年越しのシリーズものとなります。昨今の言語関係の流行りや事情から、言語のパッケージをDebianのパッケージ化を検討した方が良い場合と、あまり向かない場合がある状況とのことです。会場にHaskelの詳しい人、Perlに詳しい人がおり、様々に議論がなされました。
  
\end{frame}

  
\section{Debian Trivia Quiz}
\emtext{Debian Trivia Quiz}
\begin{frame}{Debian Trivia Quiz}

  Debian の常識、もちろん知ってますよね?
知らないなんて恥ずかしくて、知らないとは言えないあんなことやこんなこと、
みんなで確認してみましょう。

今回の出題範囲は\url{debian-devel-announce@lists.debian.org},
\url{debian-news@lists.debian.org} に投稿された
内容などからです。

\end{frame}

\subsection{問題}

%; whizzy-master ../debianmeetingresume201311.tex
% 以上の設定をしているため、このファイルで M-x whizzytex すると、whizzytexが利用できます。
%

\santaku
{Jessieが無事リリースされました。いつだったでしょうか？}
{2015/4/11}
{2015/4/18}
{2015/4/25}
{C}
{無事4/25にJessie(Debian 8)がリリースされました。Linux kernelは3.16.7、GNU Compiler Collection 4.9.2、GNOME 3.14、Apache 2.4.10、PHP 5.6.7などのバージョンのものが搭載されました。搭載されているパッケージ数は全部で43,000パッケージ以上もあります。}

\santaku
{Jessieで初めて追加されたものではないものが混ざっています。どれ？}
{Debian Games Blend}
{OpenJDK}
{androidsdk-tools}
{B}
{Jessieでの変更点はリリースノート(https://www.debian.org/releases/ jessie/amd64/release-notes/index.ja.html)にあります。Debian Gamesチームよりゲームの詰め合わせのDebian Blendが追加されたり、androidの開発ツールの一部も追加されたりしています。}

\santaku
{Debian GNU/Hurd 2015も2015/4/30にリリースされました。心臓部のGNU Machのバージョンはいくつ？}
{1.6}
{1.5}
{1.4}
{B}
{2015/4/10〜2015/4/15周辺でGNU/Hurd（upstream側)のアップデートが行われ、GNU Hurdはversion 0.6に、GNU Machは1.5になりました。未だにKVM/QEMUなどの仮想環境での動作が推奨の32bitシステムですが、多くのOSがクラウド環境で動作している世の中ですので、ホビー用途の利用でそろそろ注目されても良いかも？と思う次第です。}

\santaku
{2015年のGSoCに採択された、Debian MIPS portsについての開発内容は次のどれ？}
{多数のビルド出来ないパッケージを、ちゃんとビルドできるようにする。}
{新しいMIPS CPUへの対応}
{新しいMIPS CPU搭載製品への対応}
{A}
{Debian MIPS portsの多数のパッケージは再ビルドが出来ない、インストールに失敗するものが多数ある状況です。こちらを再ビルドできるように修正し、インストール出来るようにする事が採択されました。GSoCに限らず、いつでも本件の協力者募集中ですので、MIPS portsに興味ある方はパッケージの修正に是非ご協力くださいませー。}

\santaku
{遂にhttp.debian.netがdebian.orgのインフラに移動となりました。新しいURLはどれ？}
{http://http.debian.org/debian}
{http://httpredir.debian.org/debian}
{http://www.debian.org/}
{B}
{http.debian.netはHTTPのRedirectを活用し、ユーザに最も適切な場所にある、パッケージのミラー先を教えてくれるサービスです。遂に、実験的サービスの扱いであるdebian.netからDebian公式サービスであるdebian.orgへ移動となりました。http.debian.netをsource.listに設定している人は、httpredir.debian.orgへ修正をおねがいします。また、不具合は"mirrors" pseudo-packageでBTSするようです。参照：http://bugs.debian.org/mirrors　}

\santaku
{debianへのlibdvdcss/ZFSパッケージ搭載の件について、2015/5のBit From DPLで報告された状況は以下のどれ?}
{DPLがいろいろ議論して回っている状況}
{搭載日確定した}
{搭載を諦めた}
{A}
{2015/5/19現在、FTPMasters, SFLC, FSFと議論中とのこと。さらに近々Linux Foundationとも議論予定。現在、当初想定したよりも、実現に関する状況が複雑とのことで、本件いつ決着するか見えなくなったとの事です。}



\section{自然言語処理チーム(pkg-nlp-ja)とパッケージ}
\emtext{自然言語処理チーム(pkg-nlp-ja)とパッケージ}

\section{今後のイベント}
\emtext{今後のイベント}
\begin{frame}{今後のイベント}
\begin{itemize}
 \item 2015 OSC Hokkaido
  \url{}
 \item 関西エリアDebian勉強会
 \item 東京エリアDebian勉強会(6/20)
\end{itemize}
\end{frame}

\section{今日の宴会場所}
\emtext{今日の宴会場所}
\begin{frame}{今日の宴会場所}
未定
\end{frame}

\end{document}

;;; Local Variables: ***
;;; outline-regexp: "\\([ 	]*\\\\\\(documentstyle\\|documentclass\\|emtext\\|section\\|begin{frame}\\)\\*?[ 	]*[[{]\\|[]+\\)" ***
;;; End: ***
